\documentclass[12pt]{article}

  \usepackage[a4paper,left=20mm, right=25mm,top=25mm, bottom=25mm]{geometry}
  \usepackage[utf8]{inputenc}
  \usepackage[english]{babel}
  \usepackage[T1]{fontenc}

  %Graphics
  \usepackage{graphicx}
 
  %Tables
  \usepackage{array}
  \usepackage{tabularx}
  \usepackage{booktabs}
  \newcolumntype{K}[1]{S[table-format={#1}, table-alignment=center, table-number-alignment=center, input-decimal-markers={,}, output-decimal-marker={,}]}
  \newcolumntype{L}[1]{>{\raggedright\arraybackslash}p{#1}} % left-aligned with given width
  \newcolumntype{C}[1]{>{\centering\arraybackslash}m{#1}} % centered with given width
  \newcolumntype{R}[1]{>{\raggedleft\arraybackslash}p{#1}} % right-aligned with given width
  
  %Math formulas
  \usepackage{amsmath,amsfonts,amssymb,amsthm} %math symbols
    \newtheorem{thm}{Theorem} %[section]
    \newtheorem{defn}[thm]{Definition}
    \newtheorem{lem}[thm]{Lemma}
    \newtheorem{prop}[thm]{Proposition}
    \newtheorem{cor}[thm]{Corollary}

\usepackage{hyperref}

\usepackage{mathptmx}
\usepackage[scaled=.90]{helvet}
%\usepackage{courier}
    
 
%***********************************************************************
\begin{document}

\section{The Problem}
Let $n\in\mathbb{N}$ with $n\geq 2$. We define the index set $N=\{1, 2,\dots, n\}$ and two vectors $p = (p_1, p_2, \dots, p_n), q = (q_1, q_2, \dots, q_n)$ with $n\geq 2$ and $p_i, q_i \in \{0, 1\}$ for $i\in N$.
For $i,j \in N$ with $i<j$ let
\begin{align*}
\mathrm{one}_i^j(p) &=   \#\{ k \mid i< k < j, p_k=1 \}
\\
\overline{\mathrm{one}}_i^j(p) &=  \#\{ k \mid 1\leq k<i, p_k=1 \}
+
\#\{ k \mid j< k\leq n, p_k=1 \}
\\
\mathrm{zero}_i^j(p) 
&=   \#\{ k\mid i< k < j, p_k=0 \}
\\
\overline{\mathrm{zero}}_i^j(p) 
&=   \#\{ k\mid 1\leq k < i, p_k=0 \}
+  \#\{ k\mid j< k \leq n, p_k=0 \}
\end{align*}

\begin{defn}
Let the tuple $(n, p, q)$ be as described above.
If there exist $i, j \in N$ with $i<j$, $p_i = p_j$ and
\begin{align}
\mathrm{one}_i^{j}(p)
&> 
\overline{\mathrm{one}}_i^{j}(p),
\label{eq:condition for swap -> ones_greater}
\\
\mathrm{zero}_i^{j}(p)
&> 
\overline{\mathrm{zero}}_i^{j}(p),
\label{eq:condition for swap -> zeros_greater}
\end{align}
then the following transformation, which we call a swap, may be performed for $(i, j)$
\begin{align*}
\tau_{ij}(p) 
&= (p_1, p_2, \dots, p_{i-1}, 1-p_i, p_{i+1},\dots, p_{j-1}, 1-p_j, p_{j+1},\dots, p_n).
\end{align*}
If some $q$ can be reached from $p$ by a finite sequence of $m$ swaps, i.e.
\begin{align}
(\tau_{i_mj_m} \circ \cdots \circ \tau_{i_2j_2}\circ \tau_{i_1j_1}) (p) 
&= 
q,
\label{eq:p=q via swaps}
\end{align}
we also write $p\sim q$.
\end{defn}

\begin{defn}
Let $(n, p, q)$ be given with $\#\{ i \in N\mid p_i=1 \}$ and $\#\{ i \in N\mid p_i=0 \}$ odd.
If $p\sim q$ holds, we call $(n, p, q)$ solvable.
\label{defn:problem description}
\end{defn}
\noindent And the question of course is: given the tuple $(n, p, q)$, is $(n, p, q)$ solvable?

\section{Solution}
\begin{prop}
Let $(n, p, q)$ be as in Definition \ref{defn:problem description}. Then $n$ is even.
\begin{proof}
Since by assumption $\#\{ i \in N\mid p_i=1 \}$ and $\#\{ i \in N\mid p_i=0 \}$ are odd, $n=\#\{ i \in N\mid p_i=1 \} + \#\{ i \in N\mid p_i=0 \}$ must be even.
\end{proof}
\end{prop}


\begin{prop}
Let $(n, p, q)$ be as in Definition \ref{defn:problem description}. Then for every swap of $(i, j)$
\begin{align}
j-i
&\geq \frac{n}{2}+1,
\label{eq:j-i>n/2}
\end{align} 
i.e.\ the number of indices between $i$ and $j$ must be at least $\frac{n}{2}$.
\begin{proof}
As a necessary condition for a swap, conditions \eqref{eq:condition for swap -> ones_greater} and \eqref{eq:condition for swap -> zeros_greater} give
\begin{align*}
\mathrm{one}_i^{j}(p) + \mathrm{zero}_i^{j}(p)
&\geq
\overline{\mathrm{one}}_i^{j}(p) + \overline{\mathrm{zero}}_i^{j}(p) + 2.
\end{align*}
Bringing in the balance equations
\begin{align*}
\mathrm{one}_i^{j}(p) + \mathrm{zero}_i^{j}(p)
+
\overline{\mathrm{one}}_i^{j}(p) + \overline{\mathrm{zero}}_i^{j}(p) 
&=
n-2,
\\
\mathrm{one}_i^{j}(p) + \mathrm{zero}_i^{j}(p)
&=
j-i-1
\end{align*}
we obtain
\begin{align*}
\mathrm{one}_i^{j}(p) + \mathrm{zero}_i^{j}(p)
&\geq
n- (\mathrm{one}_i^{j}(p) + \mathrm{zero}_i^{j}(p))
\\
2(j-i-1)
&\geq
n
\\
j-i
&\geq \frac{n}{2}+1.
\end{align*}
\end{proof}
\end{prop}

\begin{defn}
Let $(n, p, q)$ be as in Definition \ref{defn:problem description}.
We define $L=\{1,\dots, \frac{n}{2}-1\}$ and $R=\{\frac{n}{2}+2,\dots, n\}$.
\end{defn}

\begin{cor}
Let $(n, p, q)$ be as in Definition \ref{defn:problem description}. Then the indices $\{\frac{n}{2}, \frac{n}{2}+1 \}$ can never swap. A swap $(i, j)$ always takes place with an index $i\in L$ and an index $j\in R$.
\label{cor:swap only from L and R}
\end{cor}

\begin{defn}
Let $(n, p, q)$ be as in Definition \ref{defn:problem description}. We call $(n, p, q)$ \emph{center-equal} if $p_{\frac{n}{2}}=q_{\frac{n}{2}}$ and $p_{\frac{n}{2}+1}=q_{\frac{n}{2}+1}$ holds.
\label{defn:center-equal}
\end{defn}

\begin{cor}
Let $(n, p, q)$ be as in Definition \ref{defn:problem description}. If $(n, p, q)$ is solvable, then it is center-equal.
\label{cor:center invariant}
\end{cor}

\begin{defn}
Let $(n, p, q)$ be as in Definition \ref{defn:problem description}. We call $i\in L$ swappable if
\begin{align}
\mathrm{one}_i^{n}(p)
&> 
\overline{\mathrm{one}}_i^{n}(p)
+
p_i-p_n
\quad\textrm{and}
\label{eq:condition i and n swappable -> ones greater}
\\
\mathrm{zero}_i^{n}(p)
&> 
\overline{\mathrm{zero}}_i^{n}(p)
-
p_i+p_n
\label{eq:condition i and n swappable -> zeros greater}
\end{align}
holds. Analogously, we call $j\in R$ swappable if
\begin{align}
\mathrm{one}_1^{j}(p)
&> 
\overline{\mathrm{one}}_1^{j}(p)
+
p_j-p_1
\quad\textrm{and}
\label{eq:condition 1 and j swappable -> ones greater}
\\
\mathrm{zero}_1^{j}(p)
&> 
\overline{\mathrm{zero}}_1^{j}(p)
-
p_j+p_1
\label{eq:condition 1 and j swappable -> zeros greater}
\end{align}
holds.
If all $\{i\in L \mid p_i\neq q_i\}$ and $\{j\in R \mid p_j\neq q_j\}$ are swappable, we call $(n, p, q)$ swappable.
\end{defn}
\noindent This means that $(i, n)$ can swap as soon as $p_n$ takes on the value of $p_i$, resp.\ $(1, j)$ can swap as soon as $p_1$ takes on the value of $p_j$.

\begin{prop}
Let $(n, p, q)$ be as in Definition \ref{defn:problem description}. If $i\in L, j\in R$ can swap, then both are swappable.
\begin{proof}
Since $(i, j)$ can swap, $p_i=p_j$. Moreover, by assumption inequalities \eqref{eq:condition for swap -> ones_greater} and \eqref{eq:condition for swap -> zeros_greater} hold.

Consider the case $j=n$. Then $p_i=p_n$ and inequalities \eqref{eq:condition i and n swappable -> ones greater} and \eqref{eq:condition i and n swappable -> zeros greater} follow, so $i$ is swappable.

Now consider $j<n$. Then
\begin{align*}
\mathrm{one}_i^{n}(p)
&=
\mathrm{one}_i^{j}(p)
+
\mathrm{one}_{j-1}^{n}(p)
\\
&\geq 
\mathrm{one}_i^{j}(p)
+
p_j+p_n
%\\
%\overline{\mathrm{one}}_i^{n}(p)
%&\leq
%\overline{\mathrm{one}}_i^{j}(p)
%-
%p_j-p_n
\\
&>
\overline{\mathrm{one}}_i^{j}(p)
+
p_j+p_n
\quad\textrm{by inequality \eqref{eq:condition for swap -> ones_greater}}
\\
&=
\overline{\mathrm{one}}_i^{j}(p)
+
p_i+p_n
\\
&\geq
\overline{\mathrm{one}}_i^{n}(p)
+
p_i+p_n
\\
&\geq
\overline{\mathrm{one}}_i^{n}(p)
+
p_i-p_n
\end{align*}
and so inequality \eqref{eq:condition i and n swappable -> ones greater} holds. Analogously, inequality \eqref{eq:condition i and n swappable -> zeros greater} follows, and thus $i$ is swappable.
By symmetry, $j$ is swappable as well.
\end{proof} 
\label{prop:if i j can swap, they are swappable}
\end{prop}

\begin{prop}
Let $(n, p, q)$ be as in Definition \ref{defn:problem description}. If $i\in L$ is not swappable, then no $i'\in L$ with $i'\geq i$ can swap. If $j\in R$ is not swappable, then no $j'\in R$ with $j'\leq j$ can swap.
\begin{proof}
By assumption
\begin{align*}
\mathrm{one}_i^{n}(p)
&\leq
\overline{\mathrm{one}}_i^{n}(p)
+
p_i-p_n.
\end{align*}
It follows that
\begin{align*}
\mathrm{one}_{i'}^{n}(p)
&\leq
\mathrm{one}_{i}^{n}(p)
-
p_{i'}
\\
&\leq
\overline{\mathrm{one}}_i^{n}(p)
+
p_i-p_n
-
p_{i'}
\\
&\leq
\overline{\mathrm{one}}_{i'}^{n}(p)
-p_i
+
p_i-p_n
-
p_{i'}
\\
&=
\overline{\mathrm{one}}_{i'}^{n}(p)
-p_n
-
p_{i'}
\\
&\leq
\overline{\mathrm{one}}_{i'}^{n}(p)
-p_n
+
p_{i'}
\end{align*}
and, with an analogous argument for the number of zeros, and by symmetry, the claim follows.
\end{proof}
\label{prop:i not swappable => i'>i not swappable}
\end{prop}

\begin{prop}
Let $(n, p, q)$ be as in Definition \ref{defn:problem description}. If $i\in L$ is swappable, then all $i'\in L$ with $i'\leq i$ are swappable. If $j\in R$ is swappable, then all $j'\in R$ with $j'\geq j$ are swappable.
\begin{proof}
By assumption
\begin{align*}
\mathrm{one}_i^{n}(p)
&>
\overline{\mathrm{one}}_i^{n}(p)
+
p_i-p_n.
\end{align*}
It follows that
\begin{align*}
\mathrm{one}_{i'}^{n}(p)
&\geq
\mathrm{one}_{i}^{n}(p)
+
p_{i}
\\
&>
\overline{\mathrm{one}}_i^{n}(p)
+
p_i-p_n
+
p_{i}
\\
&\geq
\overline{\mathrm{one}}_{i'}^{n}(p)
+ p_{i'}
+ p_i-p_n
+ p_{i}
\\
&\geq
\overline{\mathrm{one}}_{i'}^{n}(p)
+p_{i'}-p_n
\end{align*}
and, with an analogous argument for the number of zeros, and by symmetry, the claim follows.
\end{proof}
\label{prop:i swappable => i'<i swappable}
\end{prop}


\begin{lem}
Let $(n, p, q)$ be as in Definition \ref{defn:problem description}. An index $i\in N$ is swappable for $(n, p, q)$ if and only if it is swappable for every $(n, p', q)$ with $p'\sim p$.
\begin{proof}
We first show: $i\in L$ swappable $\Rightarrow$ $i$ is swappable for every $(n, p', q)$ with $p'\sim p$. It suffices to show that $i$ remains swappable after an arbitrary swap.
So consider a swap $(i', j)$ with $i'\neq i$.
Since $(i', j)$ can swap, they are also swappable by Proposition \ref{prop:if i j can swap, they are swappable}. Let $p' = \tau_{i'j}(p)$.
We distinguish 4 cases:
\begin{enumerate}

\item
$i'>i$ and $j=n$: By assumption
\begin{align}
\mathrm{one}_{i}^{n}(p')
&\geq
\mathrm{one}_{i'}^{n}(p')
+
p'_{i'}
\quad\textrm{since }j=n
\label{lem_eq:one_in geq one_i'n + p'i'}
\\
\overline{\mathrm{one}}_{i}^{n}(p)
&\leq
\overline{\mathrm{one}}_{i'}^{n}(p)
-
p_{i}
\label{lem_eq:one_bar_in(p) <= one_bar_i'n(p) - p_i}
\\
\mathrm{one}_{i'}^{n}(p)
&=
\mathrm{one}_{i'}^{n}(p')
\label{lem_eq:one_i'n(p) = one_i'n(p')}
\\
\overline{\mathrm{one}}_{i}^{n}(p)
&=
\overline{\mathrm{one}}_{i}^{n}(p')
\label{lem_eq:one_bar_in(p) = one_bar_in(p')}
\end{align}
we conclude
\begin{align*}
\mathrm{one}_{i}^{n}(p')
&\geq
\mathrm{one}_{i'}^{n}(p')
+
p'_{i'}
\quad\textrm{inequality \eqref{lem_eq:one_in geq one_i'n + p'i'}}
\\
&=
\mathrm{one}_{i'}^{n}(p)
+
p'_{i'}
\quad\textrm{equation \eqref{lem_eq:one_i'n(p) = one_i'n(p')}}
\\
&>
\overline{\mathrm{one}}_{i'}^{n}(p)
+
\underbrace{p'_{i'}+p_{i'}}_{=1}-p_n
\quad \textrm{by inequality \eqref{eq:condition i and n swappable -> ones greater}}
\\
&=
\overline{\mathrm{one}}_{i'}^{n}(p)
+
1-p_n
\\
&\geq
\overline{\mathrm{one}}_{i}^{n}(p)
+p_i
+
1-p_n
\quad\textrm{equation \eqref{lem_eq:one_bar_in(p) <= one_bar_i'n(p) - p_i}}
\\
&=
\overline{\mathrm{one}}_{i}^{n}(p')
+p_i
+
1-p_n
\quad\textrm{equation \eqref{lem_eq:one_bar_in(p) = one_bar_in(p')}}
\\
&=
\overline{\mathrm{one}}_{i}^{n}(p')
+p_i
+
p'_n
\\
&\geq
\overline{\mathrm{one}}_i^{n}(p')
+
p'_i-p'_n.
\end{align*}
Analogous estimates for the number of zeros give swappability of $i$ for $p'$.

\item
$i'>i$ and $j<n$: Then $p'_i = p'_n = p_i=p_n$. It follows that
\begin{align*}
\mathrm{one}_{i}^{n}(p')
&\geq
\mathrm{one}_{i'}^{j}(p')
+
p'_{i'}
+
p'_n
\\
&>
\overline{\mathrm{one}}_{i'}^{j}(p')
+
p'_{i'}
+
p'_n
\quad\textrm{since $(i', j)$ can swap in $p'$}
\\
&\geq
\overline{\mathrm{one}}_{i}^{n}(p')
+
p'_i
+
p'_n
+
p'_{i'}
+
p'_n
\\
&\geq
\overline{\mathrm{one}}_i^{n}(p')
+
p'_i-p'_n.
\end{align*}
With an analogous estimate for the number of zeros, $i$ remains swappable after the swap.

\item
$i'<i$ and $j<n$:
Here
\begin{align*}
\mathrm{one}_i^n(p')
&=
\mathrm{one}_i^n(p)
+ p'_j - p_j
\\
&> 
\overline{\mathrm{one}}_i^{n}(p)
 + p_i-p_n
 + p'_j - p_j
\quad \textrm{by inequality \eqref{eq:condition i and n swappable -> ones greater}}
\\
&=
\overline{\mathrm{one}}_i^n(p')
- p'_{i'} + p_{i'}
 + p_i-p_n
 + p'_j - p_j
\\
&=
\overline{\mathrm{one}}_i^n(p')
 + p_i-p_n.
\end{align*}
With an analogous estimate for the number of zeros, $i$ thus remains swappable after the swap.
\item It remains to consider the case $i'<i$ and $j=n$.
%\begin{align*}
%\overline{\mathrm{one}}_i^n(p')
%&=
%\overline{\mathrm{one}}_i^n(p)
%+ p'_{i'} - p_{i'}
%\\
%\mathrm{one}_i^{n}(p)
%&> 
%\overline{\mathrm{one}}_i^{n}(p)
%+
%p_i-p_n
%\\
%p_{i'} &= p_n
%\\
%p'_{i'} &= p'_n \neq p_n
%\\
%\overline{\mathrm{one}}_i^n(p')
%&=
%\overline{\mathrm{one}}_i^n(p)
%+p'_{i'}-p_{i'}
%\\
%p_i
%&=
%p'_i
%\end{align*}
Here
\begin{align*}
\mathrm{one}_i^n(p')
&=
\mathrm{one}_i^n(p)
\\
&> 
\overline{\mathrm{one}}_i^{n}(p)
 + p_i-p_n
\quad \textrm{by inequality \eqref{eq:condition i and n swappable -> ones greater}}
\\
&=
\overline{\mathrm{one}}_i^n(p')
- p'_{i'} + p_{i'}
 + p_i-p_n
\\
&=
\overline{\mathrm{one}}_i^n(p')
-p'_{i'} + p_{i}
\\
&=
\overline{\mathrm{one}}_i^n(p')
-p'_{n} + p'_{i}.
\end{align*}
With analogous reasoning for the number of zeros, $i$ remains swappable.
\end{enumerate}
Overall, swappability is thus preserved. It remains to show that non-swappability is also preserved. So suppose
\begin{align*}
\mathrm{one}_i^{n}(p)
&\leq
\overline{\mathrm{one}}_i^{n}(p)
+
p_i-p_n.
\end{align*}
\begin{enumerate}
\item
$i'>i$: This case cannot occur, since by Proposition \ref{prop:i not swappable => i'>i not swappable} such an $i'$ is not swappable.

\item
$i'<i$ and $j<n$:
Here
\begin{align*}
\mathrm{one}_{i}^{n}(p')
&=
\mathrm{one}_{i}^{n}(p)
+ p'_j - p_j
\\
&\leq
\overline{\mathrm{one}}_i^{n}(p)
+ p_i - p_n
+ p'_j - p_j
\\
&=
\overline{\mathrm{one}}_i^{n}(p')
-p'_{i'} + p_{i'}
+ p_i - p_n
+ p'_j - p_j
\\
&=
\overline{\mathrm{one}}_i^{n}(p')
+ p_i - p_n
\\
&=
\overline{\mathrm{one}}_i^{n}(p')
+ p'_i - p'_n
\end{align*}
and, by the same argument as before, $i$ remains non-swappable.

\item $i'<i$ and $j=n$:
Here
\begin{align*}
\mathrm{one}_{i}^{n}(p')
&=
\mathrm{one}_{i}^{n}(p)
\\
&\leq
\overline{\mathrm{one}}_i^{n}(p)
+ p_i - p_n
\\
&=
\overline{\mathrm{one}}_i^{n}(p')
-p'_{i'} + p_{i'}
+ p_i - p_n
\\
&=
\overline{\mathrm{one}}_i^{n}(p')
-p'_{i'} + p_{i} 
\\
&=
\overline{\mathrm{one}}_i^{n}(p')
- p'_n + p'_i
\end{align*}
and again, by the same argument as above, $i$ remains non-swappable.
\end{enumerate}
Overall, it follows that non-swappability of $i\in L$ is preserved under swaps, and by symmetry we can extend this property to $i\in N$.
\end{proof}

\label{lem:swappable once swappable always}
\end{lem}

\begin{cor}
Given $(n, p, q)$ swappable. Every $(n, p', q)$ with $p'\sim p$ is swappable.
\label{cor:swappable once swappable always}
\end{cor}


\begin{defn}
Let $(n, p, q)$ be as in Definition \ref{defn:problem description}. We define the left balance $b_L(p)$ and the right balance $b_R(p)$ by
\begin{align*}
b_L(p) &=\#\{i\in L \mid p_i=1, q_i=0\}
-
\#\{i\in L \mid p_i=0, q_i=1\}
\\
b_R(p) &= \#\{j\in R \mid p_j=1, q_j=0\}
-
\#\{j\in R \mid p_j=0, q_j=1\}
\end{align*}
We call $p$ balanced if
\begin{align*}
b_L(p) &= b_R(p).
\end{align*}
\end{defn}


\begin{lem}
 Let $(n, p, q)$ be as in Definition \ref{defn:problem description}. For $(n, p, q)$ to be solvable, $p$ must be balanced.
\begin{proof}
By Corollary \ref{cor:swap only from L and R}, every swap increases or decreases the left balance and the right balance simultaneously by 1. In particular, the difference between the two balances never changes. Since $b_L(p)=b_R(p)=0$ trivially holds in the case $p=q$, $p$ must be balanced.
\end{proof}
\label{lem_balanced}
\end{lem}

\begin{defn}
Let $(n, p, q)$ be as in Definition \ref{defn:problem description}. We define
\begin{align*}
\sigma(n, p, q) 
&=\#\{i\in N \mid p_i\neq q_i\}.
\end{align*}
\end{defn}

\begin{lem}
Given $(n, p, q)$ swappable and balanced with $\sigma(n, p, q)>0$. Then there exists $(n, p', q)$ swappable with $p'\sim p$ and $\sigma(n, p, q)\geq\sigma(n, p', q)$, where swappable $i\in L$ and $j\in R$ exist with $p_i=p_j\neq q_i=q_j$.
\begin{proof}
 For $i\in L$ we may assume without loss of generality that $p_i=0, q_i=1$.
 Since the problem is balanced, there is either a $j\in R$ with $p_j=0, q_j=1$ or an $i'\in L$ with $p_{i'}=1, q_{i'}=0$. In the first case we would already be done. In the second case we use the fact that $i$ and $i'$ are swappable, and can swap one of the two with $n$.
\end{proof}
\label{lem:swap on both sides}
\end{lem}


\begin{thm}
Let $(n, p, q)$ be as in Definition \ref{defn:problem description}. Then $(n, p, q)$ is solvable if and only if it is center-equal, balanced, and swappable.
\begin{proof}
That a solvable $(n,p,q)$ must be center-equal follows from Corollary \ref{cor:center invariant}. That a solvable $p$ must also be balanced follows from Lemma \ref{lem_balanced}. That a solvable $p$ also implies that $(n, p, q)$ is swappable is clear from Corollary \ref{cor:swappable once swappable always} as well. For the converse direction we assume that $(n,p,q)$ is center-equal, balanced, and swappable, and suppose the problem were nevertheless not solvable. Then we can keep performing swaps until we reach a smallest number of unsatisfiable swaps. Consider the corresponding $\bar{p}\sim p$. Then $\sigma(n, \bar{p}, q)>0$. Moreover, since $\bar{p}\sim p$, $(n, \bar{p}, q)$ is also swappable by Corollary \ref{cor:swappable once swappable always}.
Since $(n,p,q)$ is center-equal, $p_{\frac{n}{2}}=q_{\frac{n}{2}}$ and $p_{\frac{n}{2}+1}=q_{\frac{n}{2}+1}$ hold, and by Corollary \ref{cor:center invariant} applied to $\bar p\sim p$, we also have $\bar p_{\frac{n}{2}}=q_{\frac{n}{2}}$ and $\bar p_{\frac{n}{2}+1}=q_{\frac{n}{2}+1}$. Hence $\sigma(n,\bar p, q)>0$ can only be caused by an index in $L\cup R$, i.e.\ there is $i\in L$ with $\bar p_i\neq q_i$ or $j\in R$ with $\bar p_j\neq q_j$ -- exactly the hypothesis required by Lemma \ref{lem:swap on both sides}.
By Lemma \ref{lem:swap on both sides} we may therefore assume that $i\in L$ and $j\in R$ exist with $p_i=p_j\neq q_i=q_j$.
Now there are two cases:
\begin{enumerate}
 \item $\bar{p}_1 = \bar{p}_n$: If $\bar{p}_i=\bar{p}_j=\bar{p}_1=\bar{p}_n$, we set $p'=(\tau_{in}\circ\tau_{1j})(\bar{p})$, otherwise $p'=(\tau_{in}\circ\tau_{1j}\circ\tau_{1n})(\bar{p})$.
 \item $\bar{p}_1 \neq \bar{p}_n$: Here we assume without loss of generality that $\bar{p}_i=\bar{p}_n$. We then set $p'=(\tau_{1j}\circ\tau_{1n}\circ\tau_{in})(\bar{p})$.
\end{enumerate}
In every case, by performing further swaps we obtain $(n, p', q)$ with $\sigma(n, p', q)<\sigma(n, \bar{p}, q)$, a contradiction.
\end{proof}
\label{thm}
\end{thm}

\begin{cor}
It can be determined in $O(n)$ time whether the problem is solvable.
\begin{proof}
By Theorem \ref{thm} we only need to check whether $(n, p, q)$ is center-equal, balanced, and swappable.
Whether $(n,p,q)$ is center-equal can be checked directly from the two values $p_{\frac{n}{2}}, q_{\frac{n}{2}}$ and $p_{\frac{n}{2}+1}, q_{\frac{n}{2}+1}$ in $O(1)$. The balances can be computed in $O(n)$. Afterwards, by Proposition \ref{prop:i swappable => i'<i swappable}, one only needs to check whether $\max_{i\in L, p_i\neq q_i} i$ is swappable with $n$ and whether $\min_{j\in R, p_j\neq q_j} j$ is swappable with $1$. Each of these can be done in $O(n)$ time with $O(1)$ memory.
\end{proof}
\end{cor}

\end{document}
